\chapter{Contexto del Problema}

\section{Contexto Organizacional}

ProInversión (Agencia de Promoción de la Inversión Privada) es el organismo técnico
especializado adscrito al Ministerio de Economía y Finanzas del Perú, responsable
de promover, estructurar y adjudicar proyectos de inversión privada en infraestructura
y servicios públicos mediante Asociaciones Público-Privadas (APP) y Proyectos en
Activos. Su cartera activa supera los USD~50{,}000 millones en valor de proyectos,
incluyendo concesiones en infraestructura de transporte, saneamiento, energía y
telecomunicaciones que impactan directamente la calidad de vida de millones de
ciudadanos peruanos.

En este contexto, los contratos de concesión constituyen el instrumento jurídico
central que regula la relación entre el Estado y el concesionario privado durante
períodos de entre 20 y 40 años. Cada contrato define con precisión las obligaciones
técnicas, financieras, ambientales y operativas de ambas partes, los mecanismos de
renegociación, los indicadores de desempeño exigibles y los procedimientos de
resolución de disputas. Dada la magnitud económica y el horizonte temporal de estos
instrumentos, su revisión técnico-legal constituye una fase crítica del ciclo de
adjudicación: un error no detectado en una cláusula puede derivar en contingencias
fiscales de gran escala, arbitrajes internacionales costosos o retrasos que afectan
la competitividad del país como destino de inversión.

\section{Descripción del Problema}

\subsection{Problema central}

La revisión técnico-legal de contratos de concesión de infraestructura en ProInversión
enfrenta un problema central de escala y eficiencia que puede formularse como la
siguiente pregunta de investigación:

\begin{quote}
\textit{¿Cómo reducir de 320 horas a menos de 1 hora el \textit{screening} inicial de contratos de concesión APP, manteniendo trazabilidad completa y sin reemplazar el juicio experto del analista?}
\end{quote}

El punto de partida de este problema es concreto y medible. El análisis del proceso
actual, basado en entrevistas con seis especialistas de ProInversión realizadas en
octubre de 2025, indica que la revisión completa de un contrato de complejidad media
requiere en promedio \textbf{320 horas} de trabajo multidisciplinario (rango:
280--400 horas según complejidad), distribuidas entre cuatro equipos:

\begin{table}[H]
\centering
\caption{Distribución del esfuerzo de revisión manual por equipo}
\label{tab:esfuerzo_revision}
\begin{tabular}{lrr}
\toprule
\textbf{Equipo} & \textbf{Horas} & \textbf{\%} \\
\midrule
Revisión legal de cláusulas         & 120 h & 37.5\% \\
Validación técnica de especificaciones & 80 h & 25.0\% \\
Lectura y familiarización inicial   & 40 h & 12.5\% \\
Revisión económico-financiera       & 40 h & 12.5\% \\
Consolidación entre áreas           & 40 h & 12.5\% \\
\midrule
\textbf{Total}                      & \textbf{320 h} & \textbf{100\%} \\
\bottomrule
\end{tabular}
\end{table}

\begin{landscape}
\begin{figure}[H]
\centering
\resizebox{\linewidth}{!}{%
\begin{tikzpicture}[
    % ── Estilos base ──
    box/.style={
        rectangle, draw=gray!60, thick, rounded corners=3pt,
        minimum height=1.1cm, align=center, text width=3.2cm,
        font=\small
    },
    yellow box/.style={box, fill=orange!12, draw=orange!50},
    green box/.style={box, fill=green!12, draw=green!50},
    blue box/.style={box, fill=cyan!8, draw=cyan!40},
    white box/.style={box, fill=white, draw=gray!50},
    opinion box/.style={box, fill=green!18, draw=green!45, rounded corners=5pt},
    final box/.style={box, fill=green!22, draw=green!55, font=\small\bfseries},
    % ── Swimlane ──
    lane bg/.style={draw=gray!30, rounded corners=8pt, inner sep=14pt,
                    fill=#1, fill opacity=0.06, draw opacity=0.4},
    lane label/.style={font=\small\bfseries\sffamily, fill=white,
                       draw=gray!35, rounded corners=3pt,
                       inner xsep=6pt, inner ysep=3pt},
    % ── Flechas ──
    arrow/.style={-Latex, thick, color=gray!55},
    % ── Días ──
    dias/.style={font=\scriptsize, text=gray!60}
]

% ════════════════════════════════════════════════════════════
% CAPA 0  — Inicio (izquierda, fila inferior)
% ════════════════════════════════════════════════════════════
\node[yellow box] (elab) at (0, 0)
    {Elaboración del\\proyecto VIC\\\dias{\ding{168}~42 días}};
\node[white box]  (opnf) at (5, 0)
    {Opiniones no formales\\(Sector/OR/MEF)\\\dias{\ding{168}~21 días}};

% ════════════════════════════════════════════════════════════
% CAPA 1  — PROINVERSIÓN (fila central)
% ════════════════════════════════════════════════════════════
\node[blue box] (solSR)  at (12.5, 0)
    {Solicita opinión a\\Sector/OR\\\dias{\ding{168}~1 día}};
\node[blue box] (solMEF) at (19.5, 0)
    {Solicita opinión\\al MEF\\\dias{\ding{168}~1 día}};
\node[blue box] (solCGR) at (26.5, 0)
    {Solicita Informe Previo\\a la CGR\\\dias{\ding{168}~3 días}};

% ════════════════════════════════════════════════════════════
% CAPA 2  — Comité (arriba-izquierda)
% ════════════════════════════════════════════════════════════
\node[green box] (comite) at (8.5, 4.5)
    {Autoriza solicitar\\opiniones\\\dias{\ding{168}~8 días}};

% ════════════════════════════════════════════════════════════
% CAPA 2  — Sector y OR (arriba-centro)
% ════════════════════════════════════════════════════════════
\node[white box]   (reqSR) at (13.5, 3)
    {Requerimiento/\\atención de info\\\dias{\ding{168}~1 día}};
\node[opinion box] (opSR)  at (17, 3)
    {Opinión Sector/OR\\\dias{\ding{168}~21 días}};

% ════════════════════════════════════════════════════════════
% CAPA 2  — MEF (arriba-derecha)
% ════════════════════════════════════════════════════════════
\node[white box]   (reqMEF) at (20.5, 3)
    {Requerimiento/\\atención de info\\\dias{\ding{168}~1 día}};
\node[opinion box] (opMEF)  at (24, 3)
    {Opinión MEF\\\dias{\ding{168}~19 días}};

% ════════════════════════════════════════════════════════════
% CAPA 2  — CGR (arriba-derecha lejana)
% ════════════════════════════════════════════════════════════
\node[white box] (reqCGR)  at (27.5, 3)
    {Requerimiento/\\atención de info\\\dias{\ding{168}~1 día}};
\node[white box] (infPrev) at (31, 3)
    {Informe Previo\\a la VIC\\\dias{\ding{168}~21 días}};

% ════════════════════════════════════════════════════════════
% CAPA 3  — Dirección Ejecutiva (arriba-extremo derecho)
% ════════════════════════════════════════════════════════════
\node[yellow box] (aprueba) at (34.5, 4.5)
    {Aprueba DI-VIC\\\dias{\ding{168}~3 días}};
\node[yellow box] (ratifica) at (38, 4.5)
    {Ratifica DI-VIC\\\dias{\ding{168}~5 días}};

% ════════════════════════════════════════════════════════════
% Nodo final
% ════════════════════════════════════════════════════════════
\node[final box] (vic) at (41.5, 4.5) {VIC\\ratificada};

% ════════════════════════════════════════════════════════════
% Swimlane backgrounds  (on background layer)
% ════════════════════════════════════════════════════════════
\begin{scope}[on background layer]
    \node[lane bg=orange!15, fit=(reqSR)(opSR)]  (laneSR)  {};
    \node[lane bg=blue!15,   fit=(reqMEF)(opMEF)] (laneMEF) {};
    \node[lane bg=gray!15,   fit=(reqCGR)(infPrev)] (laneCGR) {};
    \node[lane bg=yellow!15, fit=(aprueba)(ratifica)] (laneDE) {};
    % Fondo PROINVERSIÓN
    \node[lane bg=cyan!8,
          fit=(solSR)(solMEF)(solCGR),
          inner xsep=20pt, inner ysep=18pt] (lanePRO) {};
\end{scope}

% ════════════════════════════════════════════════════════════
% Swimlane labels
% ════════════════════════════════════════════════════════════
\node[lane label] at (comite.north)  [above=4pt] {Comité};
\node[lane label] at (laneSR.north)  [above=4pt] {Sector y OR};
\node[lane label] at (laneMEF.north) [above=4pt] {MEF};
\node[lane label] at (laneCGR.north) [above=4pt] {CGR};
\node[lane label] at (laneDE.north)  [above=4pt] {Dirección Ejecutiva};
\node[lane label] at (lanePRO.south) [below=4pt] {PROINVERSIÓN};

% ════════════════════════════════════════════════════════════
% Flechas del flujo
% ════════════════════════════════════════════════════════════
% Inicio
\draw[arrow] (elab)  -- (opnf);
% Opiniones → Comité (sube)
\draw[arrow] (opnf.north) -- ++(0,0.8) -| (comite.south);
% Comité → Solicita Sector/OR (baja)
\draw[arrow] (comite.south) ++(0.8,0) -- ++(0,-1.2) -| (solSR.north);
% Solicita → Sector y OR (sube)
\draw[arrow] (solSR.north) ++(0.3,0) -- ++(0,0.6) -| (reqSR.south);
\draw[arrow] (reqSR) -- (opSR);
% Solicita → Solicita MEF (horizontal)
\draw[arrow] (solSR.east) -- (solMEF.west);
% Solicita MEF → MEF (sube)
\draw[arrow] (solMEF.north) ++(0.3,0) -- ++(0,0.6) -| (reqMEF.south);
\draw[arrow] (reqMEF) -- (opMEF);
% Solicita MEF → Solicita CGR (horizontal)
\draw[arrow] (solMEF.east) -- (solCGR.west);
% Solicita CGR → CGR (sube)
\draw[arrow] (solCGR.north) ++(0.3,0) -- ++(0,0.6) -| (reqCGR.south);
\draw[arrow] (reqCGR) -- (infPrev);
% Informe Previo → Aprueba
\draw[arrow] (infPrev.east) -- ++(0.8,0) |- (aprueba.west);
% Aprueba → Ratifica → VIC
\draw[arrow] (aprueba) -- (ratifica);
\draw[arrow] (ratifica) -- (vic);

\end{tikzpicture}%
}
\caption{Proceso de aprobación de la Versión Inicial del Contrato (VIC). Las fases de
\textit{Elaboración del proyecto VIC} (42 días) y \textit{Opiniones no formales} (21 días)
concentran el esfuerzo de revisión manual descrito en la Tabla~\ref{tab:esfuerzo_revision}.}
\label{fig:proceso_vic}
\end{figure}
\end{landscape}

Este volumen convierte el \textit{screening} inicial en el principal cuello de botella
del ciclo de adjudicación, impactando directamente los cronogramas de proyectos de
inversión pública de alto impacto económico y social.

\subsection{Factores que agravan el problema}

El problema central de tiempo se ve agravado por cuatro factores secundarios que
reducen adicionalmente la eficacia del proceso:

\begin{itemize}
    \item \textbf{Coordinación multidisciplinaria.} La sincronización de cuatro
    equipos especializados multiplica los esfuerzos y genera retrasos adicionales
    difíciles de planificar, al ser necesario resolver discrepancias entre criterios
    de distintas especialidades antes de consolidar observaciones.

    \item \textbf{Inconsistencias detectadas tardíamente.} Referencias cruzadas
    inválidas, contradicciones procedimentales y ambigüedades frecuentemente se
    identifican en etapas avanzadas del ciclo de adjudicación o durante la ejecución
    contractual, generando costos de corrección elevados y riesgo de arbitrajes
    internacionales.

    \item \textbf{Variabilidad analítica entre revisores.} La calidad del análisis
    depende del criterio individual de cada especialista, sin estándares uniformes
    aplicados, lo que genera inconsistencias entre contratos y riesgo de omisiones
    en revisores menos experimentados.

    \item \textbf{Trazabilidad limitada.} El proceso carece de mecanismos para
    documentar sistemáticamente la cadena de decisiones, dificultando auditorías
    posteriores y limitando la transferencia de conocimiento institucional.
\end{itemize}

Estos factores no son el problema en sí, sino las condiciones que hacen que la
solución deba ir más allá de simplemente acelerar el proceso: debe también elevar
su consistencia, cobertura y transparencia, sin prescindir del criterio experto
en la etapa de interpretación y decisión.

\subsection{Complejidad estructural del documento objeto}

La magnitud del problema de tiempo se explica en parte por la complejidad inherente
del documento que debe revisarse. El Contrato de Concesión A, caso de evaluación
de este proyecto, tiene las siguientes características:

\begin{itemize}
    \item \textbf{Extensión:} 324 páginas de contenido técnico-legal especializado.
    \item \textbf{Estructura:} 41 secciones (capítulos y anexos), con jerarquías
    de hasta tres niveles de numeración.
    \item \textbf{Cláusulas:} 448 cláusulas únicas, con hasta 60 por capítulo.
    \item \textbf{Referencias cruzadas:} Más de 200 referencias internas entre
    cláusulas y secciones, además de citas a normativa sectorial externa.
    \item \textbf{Contenido mixto:} Texto legal, especificaciones técnicas de
    ingeniería, fórmulas financieras y matrices de obligaciones.
\end{itemize}

Detectar una referencia cruzada inválida en el capítulo 3 que afecta una obligación
definida en el capítulo 18 requiere mantener en memoria activa el contenido completo
del documento, una capacidad que supera los límites prácticos de la revisión humana
sin apoyo sistemático.

\section{Antecedentes de Automatización en ProInversión}

Hasta la fecha, ProInversión ha utilizado principalmente herramientas de productividad
estándar para la gestión del proceso de revisión: procesadores de texto con funciones
básicas de búsqueda y control de cambios, hojas de cálculo para seguimiento de
observaciones y sistemas de gestión documental para almacenamiento y control de
versiones. No se han implementado sistemas automatizados específicos para la detección
de inconsistencias, validación de referencias cruzadas o análisis semántico de
coherencia contractual en documentos de esta complejidad.

Los intentos previos de digitalización se han limitado a la automatización de procesos
administrativos y flujos de aprobación, sin abordar el núcleo analítico del trabajo de
revisión técnico-legal. Esta situación contrasta con desarrollos recientes en el sector
financiero y legal internacional, donde se han comenzado a aplicar técnicas de
inteligencia artificial para análisis contractual, aunque con enfoques y dominios
diferentes al contexto específico de concesiones de infraestructura pública peruana.

\section{Estado del Arte y Trabajos Relacionados}

La aplicación de \textit{Natural Language Processing} (NLP) e inteligencia artificial
al análisis contractual ha experimentado un crecimiento significativo en los últimos
años. A continuación se presentan los antecedentes más relevantes:

\subsection{Iniciativas académicas}

El dataset \textbf{CUAD} (\textit{Contract Understanding Atticus Dataset})
\cite{hendrycks2021cuad} se consolidó como referente para revisión legal automatizada,
con anotaciones de expertos sobre 510 contratos comerciales en inglés y tareas de
localización de cláusulas específicas. Si bien establece un benchmark valioso,
sus contratos son de naturaleza comercial anglosajona, de extensión significativamente
menor que los contratos de concesión de infraestructura peruana, y no incluyen
referencias cruzadas densas ni normativa sectorial local.

\textbf{ContractNLI} \cite{koreeda2021contractnli} planteó el problema de inferencia
a nivel de documento para contratos de confidencialidad (NDA), habilitando la
detección de contradicciones y verificaciones lógicas entre secciones. Este enfoque
es conceptualmente cercano al módulo de coherencia de ContractIA, aunque aplicado
a un tipo de documento radicalmente más simple.

Revisiones sistemáticas recientes \cite{singh2025survey} confirman el crecimiento
de la clasificación y extracción en contratos como subcampo de Legal NLP, con
conjuntos de datos y tareas cada vez más estandarizadas que respaldan la madurez
tecnológica para casos de uso productivos.

\subsection*{Iniciativas industriales}

\textbf{Google Cloud Contract DocAI} \cite{google_contract_docai_2021,google_document_ai}
ofrece parsers especializados para extracción de campos y cláusulas relevantes en
contratos, con foco en flujos de \textit{Contract Lifecycle Management} (CLM). Su
diseño está optimizado para contratos comerciales estándar de menor extensión y no
incorpora análisis de coherencia interna ni validación de referencias cruzadas.

El ecosistema de \textbf{IBM watsonx} \cite{ibm_contracts_ai_aws,ibm_watsonx_contracts_2024}
promueve soluciones de Contract Management AI que integran extracción de cláusulas,
detección de riesgos y validaciones de consistencia. Sin embargo, estas soluciones
están orientadas al mercado anglosajón y no pueden incorporar la normativa sectorial
específica del Perú (Ley de APP, TUO de Contrataciones del Estado, normas del
OSITRAN, OSINERGMIN, SUNASS, entre otras).

\subsection*{Brecha identificada}

Una revisión exhaustiva de la literatura y las soluciones disponibles revela una
brecha clara: \textbf{no existe ningún sistema diseñado específicamente para la
auditoría de contratos de concesión de infraestructura bajo el marco normativo
peruano}. Los contratos de APP peruanos presentan características que los diferencian
fundamentalmente de los documentos abordados por las soluciones existentes:

\begin{itemize}
    \item Mayor complejidad estructural: hasta 41 secciones con jerarquías multinivel.
    \item Mayor extensión: 200--400 páginas frente a los 10--50 de contratos comerciales.
    \item Combinación de aspectos técnicos de ingeniería con aspectos legales y financieros.
    \item Referencias densas a normativa sectorial específica del contexto peruano.
    \item Mecanismos particulares de renegociación, riesgo compartido y resolución de
    disputas propios del modelo concesional.
\end{itemize}

Esta brecha de adaptación constituye la motivación central del presente proyecto.

\section{Stakeholders}

\begin{table}[H]
\centering
\caption{Stakeholders identificados, rol y necesidad principal}
\label{tab:stakeholders}
\begin{tabular}{p{4cm}p{3.5cm}p{6cm}}
\toprule
\textbf{Perfil} & \textbf{Rol} & \textbf{Necesidad principal} \\
\midrule
Especialista Legal      & Usuario primario   & Detección de inconsistencias jurídicas y omisiones normativas \\
Especialista Técnico    & Usuario primario   & Validación de especificaciones técnicas y plazos \\
Especialista Ambiental  & Usuario primario   & Cumplimiento de normativa ambiental sectorial \\
Analista Económico      & Usuario primario   & Consistencia de fórmulas y obligaciones financieras \\
Director de ProInversión & Usuario secundario & Reporte ejecutivo de riesgos por proyecto \\
Auditor Externo         & Usuario secundario & Evidencia trazable y verificable de cada hallazgo \\
\bottomrule
\end{tabular}
\end{table}

\section{Relevancia e Impacto Esperado}

La relevancia del problema se puede cuantificar en tres dimensiones:

\begin{itemize}
    \item \textbf{Eficiencia operacional.} Reducir el tiempo de \textit{screening}
    inicial libera capacidad de los equipos especializados para enfocarse en análisis
    de mayor valor agregado: interpretación jurídica, negociación estratégica y
    gestión de riesgos complejos que requieren genuinamente juicio experto humano.
    \item \textbf{Calidad y consistencia.} La detección sistemática y uniforme de
    inconsistencias reduce la variabilidad analítica entre revisores y disminuye la
    probabilidad de que problemas críticos escapen a la revisión, protegiendo al Estado
    de contingencias fiscales durante la ejecución contractual.
    \item \textbf{Gestión del conocimiento institucional.} La trazabilidad automática
    de hallazgos crea un repositorio estructurado de precedentes y criterios de
    revisión que facilita la capacitación de nuevos especialistas y fortalece la
    memoria institucional de ProInversión.
\end{itemize}

\section*{Usuarios y Casos de Uso}

\subsection{Perfiles de usuario}

Durante el levantamiento de requisitos (octubre de 2025) se identificaron dos
categorías de usuarios según su frecuencia e intensidad de interacción con el sistema.

\textbf{Usuarios primarios} — interactúan diariamente con el sistema como parte
central de su flujo de trabajo de revisión contractual:

\begin{table}[H]
\centering
\caption{Perfiles de usuarios primarios y su interacción con el sistema}
\label{tab:usuarios_primarios}
\begin{tabular}{p{3.5cm}p{4.5cm}p{5cm}}
\toprule
\textbf{Perfil} & \textbf{Foco de revisión} & \textbf{Necesidad principal del sistema} \\
\midrule
Especialista Legal &
Cláusulas, referencias cruzadas, coherencia jurídica &
Detección automática de referencias rotas y contradicciones procedimentales \\
Especialista Técnico &
Especificaciones de ingeniería, normas técnicas &
Validación de referencias a estándares (ISO, ASTM) y consistencia técnica \\
Especialista Ambiental &
Compromisos EIA, normativa ambiental &
Identificación de obligaciones ambientales dispersas y verificación de cumplimiento \\
Analista Económico-Financiero &
Fórmulas de reajuste, garantías, mecanismos de equilibrio &
Consistencia de fórmulas financieras y detección de ambigüedades económicas \\
Director / Gerente &
Riesgos estratégicos del contrato &
Reporte ejecutivo con hallazgos priorizados por severidad \\
\bottomrule
\end{tabular}
\end{table}

\textbf{Usuarios secundarios} — acceden de forma intermitente para propósitos
específicos de gobernanza y control:

\begin{itemize}
    \item \textbf{Auditores externos} (Contraloría, firmas independientes): requieren
    trazabilidad completa e inmutable de todas las decisiones del proceso de revisión,
    incluyendo versiones analizadas, hallazgos identificados y criterios aplicados.
    \item \textbf{Equipos de supervisión del concedente}: durante la ejecución
    contractual (15--30 años), consultan cláusulas específicas y acceden a los
    fundamentos de la revisión para mantener coherencia interpretativa.
\end{itemize}

\subsection{User stories prioritarias (MUST have)}

Las necesidades de los usuarios se expresan mediante \textit{user stories}
priorizadas con metodología MoSCoW. Las siguientes cuatro historias de prioridad
\textbf{MUST} definen el núcleo funcional del MVP (el catálogo completo se
detalla en la Sección~6):

\begin{enumerate}
    \item \textbf{US-LEG-01 — Detección de referencias cruzadas rotas.}\\
    \textit{``Como especialista legal, quiero que el sistema identifique
    automáticamente todas las referencias cruzadas inválidas o inexistentes en el
    contrato, para no tener que verificar manualmente las 400--1{,}000 citas
    potenciales.''}\\
    \textit{Criterio de aceptación clave:} detección al 100\% de referencias a
    secciones inexistentes; precisión $\geq 90\%$ en referencias a cláusulas;
    tiempo de análisis $< 1$ hora para contratos de 300 páginas.

    \item \textbf{US-LEG-02 — Detección de contradicciones procedimentales.}\\
    \textit{``Como especialista legal, quiero que el sistema identifique cláusulas
    con procedimientos o plazos contradictorios, para resolver ambigüedades antes
    de la firma y prevenir disputas durante la ejecución.''}\\
    \textit{Criterio de aceptación clave:} precisión en detección $\geq 85\%$;
    cada hallazgo incluye cita textual de ambas cláusulas conflictivas.

    \item \textbf{US-TEC-01 — Validación de referencias a normas técnicas.}\\
    \textit{``Como especialista técnico, quiero verificar que todas las normas
    técnicas citadas (ISO, ASTM, normas nacionales) estén correctamente
    referenciadas, para asegurar que las especificaciones son ejecutables con
    estándares válidos.''}\\
    \textit{Criterio de aceptación clave:} identificación de todas las citas a
    normas técnicas con sus ubicaciones exactas; alerta sobre formatos incorrectos.

    \item \textbf{US-DIR-01 — Reporte ejecutivo con hallazgos priorizados.}\\
    \textit{``Como director, quiero recibir un reporte consolidado de inconsistencias
    ordenadas por severidad, para tomar decisiones de aprobación o corrección sin
    revisar manualmente el documento completo.''}\\
    \textit{Criterio de aceptación clave:} reporte generado en formato estructurado
    con clasificación automática alta/media/baja; tiempo de generación incluido en
    el tiempo total de análisis $< 1$ hora.
\end{enumerate}